\documentclass{article}
\usepackage{tcee-repaso}
\tcetema[Cohesión, política social y convergencia en la UE]{3.B.40}{La Cohesión Económica y Social en la Unión Europea: política regional e instrumentos presupuestarios y financieros. Política social y de empleo. El proceso de convergencia real en la Unión Europea.}
\tcefecha{14/11/2026}
\begin{document}
\cabecerarepaso
\begin{columnas}

% ==========================================================================
\bloque{Esquema del cante}
\esquemacante{%
\apartadocante{Introducción: «Dios da pan a quien no tiene dientes»}{2,5}
\apartadocante{I. Política regional}{11}
\subcante{1. Importancia y justificación (integración y localización)}{4}
\subcante{2. Marco jurídico}{1}
\subcante{3. Evolución histórica}{1,5}
\subcante{4. Instrumentos 2021-2027}{3}
\subcante{5. Gestión y futuro}{1,5}
\apartadocante{II. Políticas social y de empleo}{7}
\subcante{1. Política social}{4}
\subcante{2. Política de empleo}{3}
\apartadocante{III. Convergencia real}{8,5}
\subcante{1. Convergencia $\beta$ y $\sigma$}{2,5}
\subcante{2. Evidencia: Este y Sur}{2}
\subcante{3. Papel de los fondos (discontinuidad)}{3}
\subcante{4. Dificultades de medición}{1}
\apartadocante{Conclusión}{1}}
\emph{Gráficos del cante}: mercado doméstico (\hyperref[rep:I.1]{I.1}); historia (\hyperref[rep:I.3]{I.3}); convergencia $\beta$ regional y España frente al Este y el Sur (\hyperref[rep:III]{III}); discontinuidad en el 75\,\% (\hyperref[rep:III.3]{III.3}). Fórmulas: $s_n=\frac12+\frac{1+\phi}{1-\phi}(s_E-\frac12)$; $\gamma^y=\beta_0+\beta_1\ln y+\vect{\eta}\cdot\vect{z}+u$; $g_i=\alpha+\tau D_i+f(x_i-75)+\varepsilon_i$.

% ==========================================================================
\bloque[2,5]{Introducción}
\begin{itemize}
\item \textbf{Enganche}: «Dios da pan a quien no tiene dientes». Desde 2000, Polonia 48\,\% $\to$ 81\,\% de la renta media (PPA) y Rumanía 26\,\% $\to$ 78\,\%; España 97\,\% $\to$ 92\,\%, Italia 122\,\% $\to$ 96\,\%, Grecia 85\,\% $\to$ 68\,\%. ¿Sirve la cohesión para que las regiones pobres alcancen a las ricas, o solo compensa a quien pierde con la integración?
\item \textbf{Relevancia}: mayor política de inversión (392.000 M€ en 2021-2027, $\approx$1/3 del MFP) \vertema{3.B.37}; propuesta de planes nacionales y regionales desde 2028; España: 36.700 M€, sin converger desde 2008 \vertema{4.A.17}.
\item \textbf{Contextualización}: \concepto{cohesión económica, social y territorial} (art. 174 TFUE), competencia compartida (art. 4). 2 brazos: \concepto{política regional} (FEDER, FC, FTJ) y \concepto{social y de empleo} (FSE+, mínimos, coordinación). \emph{Convergencia real} (rentas) frente a nominal (Maastricht) \vertema{3.B.43}.
\item \textbf{Problemática}: ¿por qué una política regional en un mercado único?, ¿con qué instrumentos?, ¿qué puede hacer la UE en lo social?, ¿ha habido convergencia y qué parte se debe a los fondos?
\item \textbf{Estructura}: \hyperref[rep:I]{I} regional; \hyperref[rep:II]{II} social y empleo; \hyperref[rep:III]{III} convergencia.
\end{itemize}

% ==========================================================================
\bloque[11]{I. Política regional\label{rep:I}}
\begin{ideaclave}
La integración no reparte por igual: al abaratar el comercio puede concentrar la actividad en el centro. La cohesión compensa y corrige con 392.000 M€ en 2021-2027, en gestión compartida.
\end{ideaclave}

\sub{I.1. Importancia y justificación\label{rep:I.1}}
\begin{itemize}
\item \emph{Importancia}: $\approx$1/3 del presupuesto desde 1994; 0,3-0,4\,\% del PIB de la UE; 83 de 244 regiones NUTS 2 bajo el 75\,\% (2024).
\item \emph{Política}: pagos laterales (FEDER 1975 tras 1973; Delors I 1988 con el mercado único, \autor{Padoa-Schioppa} 1987; FC 1993 con el euro); legitimidad («venganza de los lugares que no importan», \autor{Rodríguez-Pose} 2018); solidaridad (art. 3 TUE).
\item \emph{Económica}: la neoclásica promete convergencia, pero: fallos de mercado (crédito, bienes públicos, \emph{big push}); crecimiento endógeno; causación circular (\autor{Myrdal} 1957, Nobel 1974); unión monetaria \vertema{3.B.16}; y la \emph{nueva geografía económica} (\autor{Krugman} 1991, Nobel 2008) \vertema{3.B.6}.
\item \emph{Modelo} (en anexo; capital móvil, \autor{Martin} y \autor{Rogers} 1995): 2 regiones, $s_E>\frac12$; Dixit-Stiglitz; capital móvil, dueños inmóviles; $\phi=\tau^{1-\sigma}$. Con $\pi_N=\pi_S$:
\begin{formula}
\funcionclave{s_n=\frac12+\frac{1+\phi}{1-\phi}\Big(s_E-\frac12\Big)}
\end{formula}
Mercado doméstico amplificado; aglomeración total en $\phi^{\mathrm{CP}}=\frac{1-s_E}{s_E}$. Implicaciones: la integración ($\uparrow\phi$) concentra; las transferencias ($\downarrow s_E$) frenan; infraestructura entre regiones puede vaciar la periferia («doble sentido»), la interna la ayuda; dilema crecimiento-equidad (\autor{Martin} 1999). Extensiones: centro-periferia; la U de \autor{Krugman} y \autor{Venables} (1995). Evidencia: \autor{Puga} (2002): convergen países, divergen regiones.
\end{itemize}

\graficorepaso[0.95\linewidth]{mercado-domestico}{La integración concentra; la transferencia la frena}

\sub{I.2. Marco jurídico\label{rep:I.2}}
\begin{itemize}
\item \emph{Competencias}: compartida (art. 4.2.c TFUE). \emph{Marco}: arts. 174-178 TFUE (174: zonas rurales, transición industrial, desventajas permanentes; «territorial» desde Lisboa; 175: informe trienal; 176: FEDER; 177: FC). \emph{Decisión}: procedimiento ordinario con consulta al CESE y al Comité de las Regiones (1994); dinero por unanimidad en el MFP \vertema{3.B.36}.
\end{itemize}

\sub{I.3. Evolución histórica\label{rep:I.3}}
\begin{itemize}
\item 1957 FSE; 1975 FEDER por cuotas; 1986 Acta Única; \emph{1988}: duplicación y 4 principios (concentración, programación, asociación, adicionalidad), objetivo 1 $<$75\,\%; 1993 FC; 2004 Este; 2014 concentración temática; 2021 FTJ; 2025 revisión intermedia. España, mayor beneficiario 1989-2006; Polonia desde 2007.
\end{itemize}

\graficorepaso[\linewidth]{historia-cohesion}{69 $\to$ 168 $\to$ 235 $\to$ 347 $\to$ 351 $\to$ 392 miles de M€}

\sub{I.4. Instrumentos 2021-2027\label{rep:I.4}}
\begin{itemize}
\item \emph{RDC} 2021/1060 para 8 fondos; \emph{5 objetivos políticos} (inteligente, verde, conectada, social, próxima) + FTJ.
\item \emph{Fondos} (M€): FEDER 214.300 (+ Interreg 9.400); FC 39.000 (+ 11.300 al MCE; 15 Estados, RNB $<$90\,\%); FSE+ 95.100; FTJ 22.100 (19.700 hoy).
\item \emph{Regiones} (PIB PPA 2015-2017): $<$75\,\% (85\,\% UE; Andalucía, C.-La Mancha, Extremadura, Ceuta, Melilla); 75-100\,\% (60-70\,\%; 9 en España); $>$100\,\% (40-50\,\%; 5).
\item \emph{Reglas}: método de Berlín (61\,\% a menos desarrolladas); concentración (85\,\% OP1+OP2 en ricos; 40\,\% OP1 intermedios; 25\,\% pobres; 30\,\% OP2); clima 30\,\% FEDER, 37\,\% FC; condiciones favorecedoras; $n+3$.
\item \emph{España}: AA de 36.700 M€ (FEDER 23.400, FSE+ 11.300, FTJ 870, FEMPA 1.120); 35.600 en programas (724 €/hab.), 3.ª tras Polonia e Italia.
\item \emph{STEP} (2024/795): 15.200 M€. \emph{Revisión intermedia} (2025/1913 y 1914): competitividad, defensa, vivienda, agua, energía, frontera oriental; +4,5\,\% prefinanciación, +10 pp; 34.600 M€ reasignados (España 3.200).
\item \emph{Otros}: FEADER \vertema{3.B.38}; FSUE (Reserva de Solidaridad, 1.016 M€ de 2018/año); IPA III; instrumentos financieros; BEI (42.800 M€, 48\,\% en regiones de cohesión, 2025) \vertema{3.B.30}; InvestEU (garantía 26.200, 372.000 movilizados).
\end{itemize}

\graficorepaso[\linewidth]{marco-cohesion}{Un reglamento, 5 fondos, 5 objetivos y 3 categorías de regiones}

\graficorepaso[0.9\linewidth]{cohesion-estados}{En importe, Polonia, Italia y España; por habitante, el Este (España, 724 €; UE, 815)}

\sub{I.5. Gestión y futuro\label{rep:I.5}}
\begin{itemize}
\item Acuerdo de asociación (con el Semestre) $\to$ programas $\to$ autoridades de gestión y auditoría, comité de seguimiento $\to$ reembolso. Condicionalidad macro (art. 19 RDC) y del Estado de Derecho (Hungría).
\item \emph{Ejecución} (nov. 2026): pagado 30\,\% (21\,\% sin anticipos); España 20\,\% (12\,\%). 2014-2020: 99,6\,\%.
\item \emph{2028-2034}: un Fondo de 865.000 M€ con planes nacionales y regionales (27 + Interreg en vez de $\approx$540 programas); $\geq$218.000 M€ para menos desarrolladas; pagos por hitos. Críticas: centralización (Comité de las Regiones, PE, Amigos de la Cohesión). \emph{Retos}: absorción, regiones estancadas, no diluir la convergencia.
\end{itemize}

% ==========================================================================
\bloque[7]{II. Políticas social y de empleo\label{rep:II}}
\begin{ideaclave}
Poco dinero y pocas competencias: mínimos, objetivos y FSE+; los Estados se reservan salarios y protección social (art. 153.5).
\end{ideaclave}

\sub{II.1. Política social\label{rep:II.1}}
\begin{itemize}
\item \emph{Justificación}: \emph{dumping} social, movilidad, legitimidad.
\item \emph{Marco}: compartida «en los aspectos definidos» (art. 4.2.b); coordinación (5.3); cláusula horizontal (9); objetivos (151); art. 153: mínimos por mayoría cualificada (salud, condiciones, consulta, excluidos, igualdad); \emph{unanimidad} en seguridad social, despido, representación, terceros países; \emph{153.5}: no salarios, asociación, huelga ni cierre patronal; diálogo social (154-155); Comité de Protección Social (160).
\item \emph{Historia}: 1957 FSE e igualdad salarial; 1961 Carta Social; 1989 Carta comunitaria; 1992 Protocolo social; 1997 Ámsterdam; 2000 Carta de Derechos (vinculante 2009); 2010 Europa 2020.
\item \emph{Pilar} (Gotemburgo, 17-11-2017): 20 principios, 3 capítulos. Plan de 2021 (Oporto): 78\,\% de empleo, 60\,\% de adultos en formación, $-$15 M en AROPE (\emph{at risk of poverty or social exclusion}: renta $<$60\,\% de la mediana, privación material y social severa o muy baja intensidad de empleo; 5 M niños). 2025: empleo 76,1\,\% (España 72,4; objetivo nacional 76); AROPE 92,7 M (2019: 95,5); España 25,7\,\%. Renovación (julio de 2026), Estrategia contra la Pobreza (mayo de 2026), Plan de Vivienda (dic. 2025).
\item \emph{Normas}: \emph{salarios mínimos} (2022/2041: referencia 60\,\% mediana / 50\,\% media; plan si cobertura $<$80\,\%); \emph{TJUE C-19/23} (11-11-2025): anula art. 5.2 y la cláusula de 5.3, el resto vale. Transparencia salarial (2023/970, 7-6-2026); plataformas (2024/2831, presunción de laboralidad, 2-12-2026; ley rider 12/2021).
\item \emph{FSE+}: $\approx$95.000 M€ (87.995 de 2018); integra FEAD; reservas 25\,\% inclusión, 3\,\% privación, 12,5\,\% jóvenes, 5\,\% pobreza infantil.
\end{itemize}

\graficorepaso[0.95\linewidth]{pilar-social}{El empleo se acerca al 78\,\%; la pobreza apenas baja}

\sub{II.2. Política de empleo\label{rep:II.2}}
\begin{itemize}
\item \emph{Marco}: solo coordinación (5.2); arts. 145-150; orientaciones anuales y recomendaciones (148); sin armonización (149); Comité de Empleo (150).
\item \emph{EEE} (Luxemburgo 1997, MAC); flexiseguridad; en el \emph{Semestre} desde 2011: informe conjunto, cuadro social, marco de convergencia social (España entre los 9 de 2026), orientaciones (Decisión 2025/2254) \vertema{3.B.45}.
\item \emph{Garantía Juvenil} (2013, reforzada 2020; ninis 11\,\% en 2025, objetivo 9\,\%).
\item \emph{FEAG} (2021/691: $>$200 despidos; 30 M€ de 2018/año); EURES; \emph{SURE}: 98.400 M€ a 19 Estados; España 21.324; 31,5 M trabajadores en 2020 \vertema{4.A.16}.
\end{itemize}

% ==========================================================================
\bloque[8,5]{III. Convergencia real\label{rep:III}}
\begin{ideaclave}
Convergencia rápida, pero la explica el Este; el Sur se aleja desde 2008. Los fondos ayudan cerca del 75\,\% si hay capacidad de absorción.
\end{ideaclave}

\sub{III.1. Convergencia $\beta$ y $\sigma$\label{rep:III.1}}
\begin{itemize}
\item \autor{Sala-i-Martín} (1990), \autor{Barro} y \autor{Sala-i-Martín} (1991) \vertema{3.A.45}:
\begin{formula}
\funcionclave{\gamma^{y}_{t,t-1}=\beta_0+\beta_1\ln y_{t-1}+\vect{\eta}\cdot\vect{z}_{t-1}+u_t}
\end{formula}
absoluta ($\vect\eta=\vect0$) o condicional; $\beta_1=-(1-e^{-\lambda T})/T$; $\sigma^2_{y,t-1}>\sigma^2_{y,t}$; $\beta$ necesaria, no suficiente. Clubes (\autor{Quah} 1996); trampa del desarrollo (\autor{Diemer} et al. 2022): 28\,\% de la población; trampa del talento (46 regiones).
\item \emph{Estimación propia} (232 NUTS 2, 2000-2024): $\hat\beta_1=-1{,}99$ (0,12), $R^2=0{,}55$, $\lambda\simeq2{,}7$\,\%; sin el Este, pendiente nula. $\sigma$: Estados 0,56 $\to$ 0,31; regiones 0,46 $\to$ 0,34 (parada 2008-2016). Clásico: 2\,\%.
\end{itemize}

\graficorepaso[0.9\linewidth]{sigma-regional}{La dispersión baja; entre regiones se para en 2008-2016}

\graficorepaso[0.95\linewidth]{convergencia-regional}{Convergencia $\beta$ entre regiones: el Este tira, el resto no}

\sub{III.2. Evidencia: el Este y el Sur\label{rep:III.2}}
\begin{itemize}
\item Hasta 2008, convergencia (top/bottom 20\,\%: 3,5 $\to$ 2,8); 2008-2013, la crisis (27 regiones con paro $>$20\,\%); después, 2 velocidades: Este 45\,\% (1995) $\to\approx$80\,\%; Chequia = España (92); Italia 96, Grecia 68, Portugal 81. España: 91 (1995) $\to$ 105 (2006) $\to$ 82 (2020) $\to$ 92 (2025) \vertema{4.A.28}.
\end{itemize}

\graficorepaso[\linewidth]{pib-pc-paises}{España, entre un Sur que se aleja y un Este que llega}

\sub{III.3. Papel de los fondos\label{rep:III.3}}
\begin{itemize}
\item \emph{Discontinuidad} (\autor{Becker}, \autor{Egger} y \autor{von Ehrlich} 2010): $D_i=\mathbb 1[x_i<75]$, sin manipulación, continuidad:
\begin{formula}
\funcionclave{g_i=\alpha+\tau D_i+f(x_i-75)+\varepsilon_i}
\end{formula}
$\tau$ = salto (versión borrosa). Más PIB por habitante ($\approx$1,6 pp más de crecimiento anual en el periodo), no más empleo; coste-eficiente. Extensiones: absorción (2013: solo $\approx$30\,\% de regiones), persistencia (2018: se deshace), \autor{Pellegrini} et al. (2013). Evidencia: QUEST/RHOMOLO (1 € $\to$ 1,3 € de PIB en 2030), HERMIN (5-10\,\%), \autor{De la Fuente} y \autor{Vives} (1995), \autor{Boldrin} y \autor{Canova} (2001). Valoración: efecto local.
\end{itemize}

\graficorepaso[0.9\linewidth]{rdd-umbral}{El salto en el 75\,\% es el efecto}

\sub{III.4. Dificultades de medición\label{rep:III.4}}
\begin{itemize}
\item Variables omitidas; contrafactual (expulsión frente a atracción; adicionalidad); efectos fuera (contratos del centro); PPA y conmutación; en la balanza de pagos: estructurales en capital, FSE+ en renta secundaria \vertema{3.B.11}.
\end{itemize}

% ==========================================================================
\bloque[1]{Conclusión}
\begin{itemize}
\item \textbf{Recapitulación}: la integración concentra; la cohesión compensa y corrige; en lo social, mínimos y coordinación; convergencia a 2 velocidades.
\item \textbf{Relevancia}: para España, no cuánto recibe, sino en qué lo gasta \vertema{4.A.17} \vertema{4.A.30}.
\item \textbf{Extensiones}: \vertema{3.B.37} \vertema{3.B.6} \vertema{3.A.45} \vertema{4.B.14} \vertema{4.A.16}.
\item \textbf{Opinión}: invertir en capacidades, no compensar ni hacer de caja; planes nacionales sin dejar fuera a las regiones.
\item \textbf{Idea final}: el pan sí alimenta, pero solo a quien puede masticarlo: hay que ayudar a las regiones a tener dientes.
\end{itemize}

% ==========================================================================
\bloque{Datos clave}
\begin{itemize}
\item \emph{Dinero}: cohesión 2021-2027, 392.000 M€ ($\approx$1/3 del MFP); en programas, 367.000 sin Interreg; Polonia 75.500, Italia 42.200, España 35.600. España: 23.400 FEDER, 11.300 FSE+, 870 FTJ.
\item \emph{Historia}: 69.000 M€ (1989-1993, 25\,\% del presupuesto); 168.000 (1994-1999); 235.000 (2000-2006); 347.000 (2007-2013); 351.000 (2014-2020). España recibió 14.200, 42.400 y 56.300 M€ en 1989-2006.
\item \emph{Revisión intermedia}: 34.600 M€ reasignados en 186 programas: competitividad 15.200, defensa 11.900, vivienda 3.300, agua 3.100, energía 1.200.
\item \emph{Regiones españolas} (\% de la media, 2000 $\to$ 2024): Madrid 129 $\to$ 125; Andalucía 72 $\to$ 69; Extremadura 64 $\to$ 70; Canarias 91 $\to$ 72; C. Valenciana 93 $\to$ 77. Capitales del Este: Varsovia 105 $\to$ 157.
\item \emph{Social} (2025): empleo UE 76,1\,\%, España 72,4\,\%; AROPE UE 20,9\,\% (92,7 M), España 25,7\,\% (niños, 34\,\%); ninis 11\,\%.
\item \emph{Absorción} (nov. 2026): 30\,\% pagado (España 20\,\%); 2014-2020: 99,6\,\%.
\end{itemize}

% ==========================================================================
\bloque{Autores y fechas}
\begin{itemize}
\item 1957 Roma (FSE); 1975 FEDER; 1987 \autor{Padoa-Schioppa}; 1988 Delors I; 1991 \autor{Barro} y \autor{Sala-i-Martín}, \autor{Krugman}; 1993 FC; 1995 \autor{Martin} y \autor{Rogers}, \autor{Krugman} y \autor{Venables}, \autor{De la Fuente} y \autor{Vives}; 1996 \autor{Quah}; 1997 EEE; 1999 \autor{Martin}; 2001 \autor{Boldrin} y \autor{Canova}; 2002 \autor{Puga}; 2010 \autor{Becker} et al.; 2017 Pilar; 2018 \autor{Rodríguez-Pose}; 2021 RDC; 2022 salarios mínimos; 2024 9.º informe; 2025 revisión intermedia y TJUE; 2026 renovación del Pilar.
\end{itemize}

\bloque{Preguntas del tribunal}
\begin{enumerate}
\item \textbf{¿Por qué cohesión si la neoclásica predice convergencia?} Fallos de mercado, rendimientos crecientes, NGE: la integración concentra.
\item \textbf{¿Toda infraestructura ayuda?} No: la que une con el centro puede vaciar la periferia.
\item \textbf{¿Qué es la adicionalidad?} Los fondos se suman a la inversión nacional, no la sustituyen.
\item \textbf{¿Hay un salario mínimo europeo?} No: procedimiento y referencias; el TJUE anuló los criterios obligatorios (153.5).
\item \textbf{¿$\beta$ implica $\sigma$?} No (adelantamiento, perturbaciones); al revés, sí.
\item \textbf{¿Funcionan los fondos?} Sí cerca del 75\,\% y con capacidad de absorción; efecto local y poco persistente.
\item \textbf{¿Qué cambia en 2028?} Un fondo con planes nacionales y regionales y pagos por hitos.
\item \textbf{¿Por qué converge el Este y no el Sur?} Cadenas de valor e inversión extranjera, capital humano; el Sur: competencia de Asia, crisis del euro, baja productividad, poca absorción.
\item \textbf{¿Cómo se registran los fondos en la balanza de pagos?} FEDER y FC, transferencias de capital; FSE+, renta secundaria.
\item \textbf{¿FEAG o FSE+?} FEAG: reacción puntual a despidos masivos, fuera del MFP; FSE+: programación a 7 años.
\end{enumerate}

\begin{trampa}
\begin{itemize}
\item Regiones: \emph{75\,\%} y \emph{100\,\%} del PIB por hab. en PPA (en 2014-2020, 75 y \emph{90\,\%}); Fondo de Cohesión: \emph{RNB} $<$\emph{90\,\%} y por \emph{Estado}, no por región.
\item Cofinanciación UE: \emph{85\,\%} / \emph{60\,\%} / \emph{40\,\%}. Liberación: \emph{$n+3$}.
\item Empleo: competencia de \emph{coordinación}; social: \emph{compartida} en lo definido. Seguridad social: \emph{unanimidad}, sin pasarela.
\item FSE (\emph{1957}) es el fondo más antiguo; FEDER \emph{1975}; FC \emph{1993}.
\item Objetivo 2030 AROPE: $-$\emph{15} millones frente a \emph{2019}; empleo \emph{78\,\%} (20-64 años); formación \emph{60\,\%}.
\item Directiva de salarios mínimos: \emph{no} fija un salario europeo; umbral de negociación colectiva del \emph{80\,\%}.
\item Pilar: \emph{20} principios, \emph{3} capítulos, Gotemburgo \emph{2017}; no vinculante.
\end{itemize}
\end{trampa}

\end{columnas}
\end{document}
